\chapter{电磁场的四维协变形式}

\section{本章目标与符号约定}

前面几章已经建立了本章所需的基本语言：闵可夫斯基时空、洛伦兹变换、四维张量、相对论粒子力学以及经典场论的拉格朗日形式。本章在这些基础上讨论电磁场的四维协变形式。

本章采用四维张量语言重写经典电动力学。电磁相互作用由四维势描述，电场和磁场合并为反对称二阶张量，Maxwell 方程则表现为洛伦兹协变的张量方程。

本章采用闵可夫斯基度规号差约定
\begin{equation}
	\eta_{\mu\nu}=\operatorname{diag}(1,-1,-1,-1).
	\label{eq:8.1}
\end{equation}
四维坐标取为

\[
x^{\mu}=(ct,x,y,z).
\]

四维偏导算符为
\begin{equation}
	\partial_{\mu}
	=
	\frac{\partial}{\partial x^{\mu}}
	=
	\left(
	\frac{1}{c}\frac{\partial}{\partial t},
	\bm{\nabla}
	\right),
	\label{eq:8.3}
\end{equation}
升高指标后有

\[
\partial^{\mu}
=
\eta^{\mu\nu}\partial_{\nu}
=
\left(
\frac{1}{c}\frac{\partial}{\partial t},
-\bm{\nabla}
\right).
\]

四维电流取为
\begin{equation}
	j^{\mu}=(c\rho,\bm{j}),
	\qquad
	j_{\mu}=(c\rho,-\bm{j}).
	\label{eq:8.5}
\end{equation}
电荷守恒的四维形式为
\begin{equation}
	\partial_{\mu}j^{\mu}=0.
	\label{eq:8.6}
\end{equation}
本章采用 SI 制单位。达朗贝尔算子记为
\begin{equation}
	\dalembert
	=
	\partial_{\mu}\partial^{\mu}
	=
	\frac{1}{c^{2}}\frac{\partial^{2}}{\partial t^{2}}
	-
	\bm{\nabla}^{2}.
	\label{eq:8.7}
\end{equation}

\section{四维势与三维势}

电磁场可以用标量势 $\phi$ 和矢量势 $\bm{A}$ 描述。取
\begin{equation}
	\bm{E}
	=
	-\bm{\nabla}\phi
	-
	\frac{\partial\bm{A}}{\partial t},
	\qquad
	\bm{B}
	=
	\bm{\nabla}\times\bm{A}.
	\label{eq:8.8}
\end{equation}
在相对论形式中，标量势和矢量势不是两个独立对象，而是同一个四维势的不同分量。本章采用
\begin{equation}
	A^{\mu}
	=
	\left(
	\frac{\phi}{c},
	\bm{A}
	\right).
	\label{eq:8.9}
\end{equation}
由度规 \eqref{eq:8.1} 降低指标，得
\begin{equation}
	A_{\mu}
	=
	\eta_{\mu\nu}A^{\nu}
	=
	\left(
	\frac{\phi}{c},
	-\bm{A}
	\right).
	\label{eq:8.10}
\end{equation}
因此
\[
A_{0}=\frac{\phi}{c},
\qquad
A_{i}=-A^{(3)}_{i},
\]
其中 $A^{(3)}_{i}$ 表示三维矢量势的第 $i$ 个分量。后文仍用 $\bm{A}$ 表示三维矢量势。

由于 $A^{\mu}$ 是四维矢量，在洛伦兹变换下满足
\begin{equation}
	A'^{\mu}=\Lambda^{\mu}{}_{\nu}A^{\nu}.
	\label{eq:8.12}
\end{equation}
例如，设参考系 $K'$ 相对于 $K$ 沿 $x$ 轴正方向以速度 $V$ 运动，从 $K$ 到 $K'$ 的 boost 给出
\[
A'^{0}=\gamma\left(A^{0}-\frac{V}{c}A^{1}\right),
\qquad
A'^{1}=\gamma\left(A^{1}-\frac{V}{c}A^{0}\right),
\]
以及
\[
A'^{2}=A^{2},
\qquad
A'^{3}=A^{3}.
\]
代入 $A^{0}=\phi/c$ 与 $A^{1}=A_{x}$，得到
\[
\phi'
=
\gamma(\phi-VA_{x}),
\qquad
A'_{x}
=
\gamma\left(A_{x}-\frac{V}{c^{2}}\phi\right),
\]
并且
\[
A'_{y}=A_{y},
\qquad
A'_{z}=A_{z}.
\]
由此可见，标量势和矢量势在洛伦兹变换下会相互混合；完整的相对论对象是四维势 $A^{\mu}$。

\section{电磁场张量}

\subsection{由四维势定义场强张量}

电磁场张量定义为
\begin{equation}
	F_{\mu\nu}
	=
	\partial_{\mu}A_{\nu}-\partial_{\nu}A_{\mu}.
	\label{eq:8.17}
\end{equation}
由定义立刻得到
\begin{equation}
	F_{\mu\nu}=-F_{\nu\mu}.
	\label{eq:8.18}
\end{equation}
因此在四维时空中，$F_{\mu\nu}$ 只有六个独立分量。这六个分量正好对应三维电场和磁场的六个分量。

若四维势按 \eqref{eq:8.12} 作为四维矢量变换，则 $F_{\mu\nu}$ 按二阶协变张量变换：
\begin{equation}
	F'_{\mu\nu}
	=
	\Lambda_{\mu}{}^{\alpha}\Lambda_{\nu}{}^{\beta}F_{\alpha\beta}.
	\label{eq:8.19}
\end{equation}
因此，电场和磁场在不同惯性系中可以相互混合，而 $F_{\mu\nu}$ 作为整体具有洛伦兹张量的变换性质。

\subsection{电场和磁场作为场张量的分量}

下面把 $F_{\mu\nu}$ 写成 $\bm{E}$ 和 $\bm{B}$ 的形式。由 \eqref{eq:8.10} 与 \eqref{eq:8.3} 可得

\[
F_{0i}
=
\partial_{0}A_{i}-\partial_{i}A_{0}.
\]

利用 $A_{i}=-A^{(3)}_{i}$，有
\begin{equation}
	F_{0i}
	=
	-\frac{1}{c}\frac{\partial A^{(3)}_{i}}{\partial t}
	-
	\frac{1}{c}\partial_{i}\phi
	=
	\frac{E_{i}}{c}.
	\label{eq:8.21}
\end{equation}
空间分量为
\begin{equation}
	F_{ij}
	=
	\partial_{i}A_{j}-\partial_{j}A_{i}
	=
	-\varepsilon_{ijk}B_{k}.
	\label{eq:8.22}
\end{equation}
因此
\begin{equation}
	(F_{\mu\nu})
	=
	\begin{pmatrix}
		0 & E_{x}/c & E_{y}/c & E_{z}/c\\
		-E_{x}/c & 0 & -B_{z} & B_{y}\\
		-E_{y}/c & B_{z} & 0 & -B_{x}\\
		-E_{z}/c & -B_{y} & B_{x} & 0
	\end{pmatrix}.
	\label{eq:8.23}
\end{equation}

\subsection{场张量的升降指标}

利用闵可夫斯基度规升高两个指标：
\begin{equation}
	F^{\mu\nu}
	=
	\eta^{\mu\alpha}\eta^{\nu\beta}F_{\alpha\beta}.
	\label{eq:8.24}
\end{equation}
由 \eqref{eq:8.23} 得
\begin{equation}
	(F^{\mu\nu})
	=
	\begin{pmatrix}
		0 & -E_{x}/c & -E_{y}/c & -E_{z}/c\\
		E_{x}/c & 0 & -B_{z} & B_{y}\\
		E_{y}/c & B_{z} & 0 & -B_{x}\\
		E_{z}/c & -B_{y} & B_{x} & 0
	\end{pmatrix}.
	\label{eq:8.25}
\end{equation}
特别地，
\begin{equation}
	F^{0i}=-\frac{E_{i}}{c},
	\qquad
	F^{i0}=\frac{E_{i}}{c},
	\qquad
	F^{ij}=F_{ij}=-\varepsilon_{ijk}B_{k}.
	\label{eq:8.26}
\end{equation}
这些符号关系将在从四维 Maxwell 方程恢复三维形式时使用。

\section{规范变换与规范不变性}

\subsection{势函数的不唯一性}

三维势 $\phi,\bm{A}$ 并不是唯一的。电场和磁场由 \eqref{eq:8.8} 给出，只要势的改变不改变 $\bm{E}$ 和 $\bm{B}$，就描述同一个物理电磁场。这种势函数的冗余称为规范自由度。

在四维形式中，这种自由度表现为四维势的一类变换。设 $f(x)$ 是充分光滑的标量函数，定义
\begin{equation}
	A'_{\mu}=A_{\mu}-\partial_{\mu}f.
	\label{eq:8.27}
\end{equation}
这称为电磁规范变换。

\subsection{四维规范变换}

在规范变换 \eqref{eq:8.27} 下，场张量变为
\[
F'_{\mu\nu}
=
\partial_{\mu}A'_{\nu}-\partial_{\nu}A'_{\mu}.
\]
代入 \eqref{eq:8.27}，得
\begin{align}
	F'_{\mu\nu}
	&=
	\partial_{\mu}(A_{\nu}-\partial_{\nu}f)
	-
	\partial_{\nu}(A_{\mu}-\partial_{\mu}f) \\
	&=
	F_{\mu\nu}
	-
	(\partial_{\mu}\partial_{\nu}f-
	\partial_{\nu}\partial_{\mu}f).
\end{align}
若 $f$ 足够光滑，混合偏导可交换，因此
\begin{equation}
	F'_{\mu\nu}=F_{\mu\nu}.
	\label{eq:8.29}
\end{equation}
规范变换改变势函数的代表元，但保持场张量不变。

\subsection{三维规范变换}

由 \eqref{eq:8.10} 与 \eqref{eq:8.27} 可得三维形式：
\begin{equation}
	\phi'=
	\phi-
	\frac{\partial f}{\partial t},
	\qquad
	\bm{A}'=
	\bm{A}+\bm{\nabla}f.
	\label{eq:8.30}
\end{equation}
代入 \eqref{eq:8.8}，电场变为
\begin{align}
	\bm{E}'
	&=
	-\bm{\nabla}\phi'
	-
	\frac{\partial\bm{A}'}{\partial t} \\
	&=
	-\bm{\nabla}\left(\phi-\frac{\partial f}{\partial t}\right)
	-
	\frac{\partial}{\partial t}(\bm{A}+\bm{\nabla}f) \\
	&=
	-\bm{\nabla}\phi
	-
	\frac{\partial\bm{A}}{\partial t}
	=
	\bm{E}.
\end{align}
磁场变为
\[
\bm{B}'
=
\bm{\nabla}\times\bm{A}'
=
\bm{\nabla}\times(\bm{A}+\bm{\nabla}f)
=
\bm{B}.
\]
由此得到三维电场和磁场的规范不变性。

\subsection{场张量的规范不变性}

式 \eqref{eq:8.29} 表明，真正由势导出的可观测对象不是单个势函数，而是规范等价类所确定的场张量 $F_{\mu\nu}$。在经典电动力学中，$A_{\mu}$ 是描述电磁相互作用的基本势变量，而 $F_{\mu\nu}$ 是规范不变的场强变量。

从带电粒子的作用量角度看，相互作用项
\[
S_{\mathrm{int}}=-q\int A_{\mu}\dif x^{\mu}
\]
在规范变换下只改变端点项：
\[
S'_{\mathrm{int}}
=
S_{\mathrm{int}}+qf(b)-qf(a).
\]
因此，在端点固定的变分问题中，规范变换不改变运动方程。

\section{Maxwell 方程的协变形式}

\subsection{齐次 Maxwell 方程}

由场张量定义 \eqref{eq:8.17} 可得恒等式
\begin{equation}
	\partial_{\lambda}F_{\mu\nu}
	+
	\partial_{\mu}F_{\nu\lambda}
	+
	\partial_{\nu}F_{\lambda\mu}
	=
	0.
	\label{eq:8.34}
\end{equation}
这是因为每一项展开后均由二阶偏导数组成，而二阶偏导可交换。式 \eqref{eq:8.34} 称为齐次 Maxwell 方程的四维协变形式。

下面通过分量展开说明这一恒等式所包含的三维内容。

\subsection{非齐次 Maxwell 方程}

非齐次 Maxwell 方程写为
\begin{equation}
	\partial_{\mu}F^{\mu\nu}
	=
	\mu_{0}j^{\nu}.
	\label{eq:8.36}
\end{equation}
这是电磁场的动力学方程，其作用量来源将在后文给出。这里先讨论它的各个分量。

\subsection{三维 Maxwell 方程}

先考察齐次方程 \eqref{eq:8.34}。当三个指标均为空间指标时，只需取
$(\lambda,\mu,\nu)=(1,2,3)$，便有
\[
\partial_{1}F_{23}
+
\partial_{2}F_{31}
+
\partial_{3}F_{12}
=0.
\]
由 \eqref{eq:8.22}，
\[
F_{23}=-B_{1},
\qquad
F_{31}=-B_{2},
\qquad
F_{12}=-B_{3},
\]
所以
\[
\partial_{1}B_{1}
+
\partial_{2}B_{2}
+
\partial_{3}B_{3}
=0,
\]
即
\[
\bm{\nabla}\cdot\bm{B}=0.
\]

再令齐次方程的一个指标为时间指标，取
$(\lambda,\mu,\nu)=(0,i,j)$，得到
\[
\partial_{0}F_{ij}
+
\partial_{i}F_{j0}
+
\partial_{j}F_{0i}
=0.
\]
利用 $F_{ij}=-\varepsilon_{ijk}B_{k}$、$F_{0i}=E_{i}/c$ 和
$\partial_{0}=c^{-1}\partial/\partial t$，上式化为
\[
-\varepsilon_{ijk}\frac{\partial B_{k}}{\partial t}
-\partial_{i}E_{j}
+\partial_{j}E_{i}
=0.
\]
例如，取 $(i,j)=(1,2)$，可得
\[
\partial_{1}E_{2}-\partial_{2}E_{1}
=
-\frac{\partial B_{3}}{\partial t}.
\]
这是旋度方程的第三个分量；另两个分量由指标的循环置换得到。因此
\[
\bm{\nabla}\times\bm{E}
=
-\frac{\partial\bm{B}}{\partial t}.
\]
由此可见，磁场的 Gauss 定律和 Faraday 电磁感应定律都直接来自
$F_{\mu\nu}$ 的定义。

下面考察非齐次方程 \eqref{eq:8.36}。先取 $\nu=0$。由
\eqref{eq:8.26} 得 $F^{i0}=E_{i}/c$，于是
\[
\partial_{\mu}F^{\mu0}
=
\partial_{i}F^{i0}
=
\frac{1}{c}\bm{\nabla}\cdot\bm{E}.
\]
而右侧为 $\mu_{0}j^{0}=\mu_{0}c\rho$。代入 \eqref{eq:8.36} 并利用 $\mu_{0}c^{2}=1/\varepsilon_{0}$，得到
\[
\bm{\nabla}\cdot\bm{E}
=
\frac{\rho}{\varepsilon_{0}}.
\]
这就是 Gauss 定律。

再取 $\nu=i$。由 \eqref{eq:8.26} 可得
\[
\partial_{\mu}F^{\mu i}
=
(\bm{\nabla}\times\bm{B})_{i}
-
\frac{1}{c^{2}}\frac{\partial E_{i}}{\partial t}.
\]
因此 \eqref{eq:8.36} 给出
\[
\bm{\nabla}\times\bm{B}
=
\mu_{0}\bm{j}
+
\frac{1}{c^{2}}\frac{\partial\bm{E}}{\partial t}.
\]
这就是 Ampere--Maxwell 方程。

将以上结果合在一起，便得到经典 Maxwell 方程组。含有电荷密度和电流密度的两式为
\begin{equation}
	\bm{\nabla}\cdot\bm{E}
	=
	\frac{\rho}{\varepsilon_{0}},
	\qquad
	\bm{\nabla}\times\bm{B}
	=
	\mu_{0}\bm{j}
	+
	\frac{1}{c^{2}}\frac{\partial\bm{E}}{\partial t},
	\label{eq:8.41}
\end{equation}
它们分别是电场的 Gauss 定律和 Ampere--Maxwell 方程。其余两式为
\begin{equation}
	\bm{\nabla}\cdot\bm{B}=0,
	\qquad
	\bm{\nabla}\times\bm{E}
	=
	-
	\frac{\partial\bm{B}}{\partial t}.
	\label{eq:8.42}
\end{equation}
它们分别是磁场的 Gauss 定律和 Faraday 电磁感应定律。这样，经典
Maxwell 方程组便作为协变方程 \eqref{eq:8.34} 与 \eqref{eq:8.36}
的三维分量形式自然出现。

\section{电荷守恒}

非齐次 Maxwell 方程 \eqref{eq:8.36} 与场张量反对称性直接推出电荷守恒。对 \eqref{eq:8.36} 两边作用 $\partial_{\nu}$，得

\[
\partial_{\nu}\partial_{\mu}F^{\mu\nu}
=
\mu_{0}\partial_{\nu}j^{\nu}.
\]

左侧为零。原因是 $\partial_{\nu}\partial_{\mu}$ 对 $\mu,\nu$ 对称，而 $F^{\mu\nu}$ 对 $\mu,\nu$ 反对称。因此
\begin{equation}
	\partial_{\mu}j^{\mu}=0.
	\label{eq:8.44}
\end{equation}
展开为三维形式，由 \eqref{eq:8.5} 得
\begin{equation}
	\frac{\partial\rho}{\partial t}
	+
	\bm{\nabla}\cdot\bm{j}=0.
	\label{eq:8.45}
\end{equation}
所以在协变电磁理论中，电荷守恒不是独立附加条件，而是非齐次 Maxwell 方程自洽性的必要结果。

\section{电磁场的作用量与场方程}

\subsection{电磁场拉格朗日密度}

按照第七章的约定，场作用量写成
\begin{equation}
	S=
	\frac{1}{c}
	\int
	\mathcal{L}\,\dif^{4}x.
	\label{eq:8.46}
\end{equation}
在 SI 制下，电磁场与给定外电流的拉格朗日密度可取为
\begin{equation}
	\mathcal{L}
	=
	-
	\frac{1}{4\mu_{0}}F_{\mu\nu}F^{\mu\nu}
	-
	j^{\mu}A_{\mu}.
	\label{eq:8.47}
\end{equation}
其中第一项是自由电磁场拉格朗日密度，第二项描述电流与电磁势的相互作用。由于 $F_{\mu\nu}F^{\mu\nu}$ 和 $j^{\mu}A_{\mu}$ 都是洛伦兹标量，作用量具有洛伦兹不变性。

\subsection{对四维势的变分}

对 $A_{\nu}$ 作变分。由 \eqref{eq:8.17} 得
\begin{equation}
	\delta F_{\mu\nu}
	=
	\partial_{\mu}\delta A_{\nu}
	-
	\partial_{\nu}\delta A_{\mu}.
	\label{eq:8.48}
\end{equation}
自由场部分的变分为

\[
\delta\mathcal{L}_{\mathrm{f}}
=
-
\frac{1}{2\mu_{0}}F^{\mu\nu}\delta F_{\mu\nu}.
\]

代入 \eqref{eq:8.48}，并利用 $F^{\mu\nu}$ 的反对称性，得到
\[
\delta\mathcal{L}_{\mathrm{f}}
=
-
\frac{1}{\mu_{0}}F^{\mu\nu}\partial_{\mu}\delta A_{\nu}.
\]
对其分部积分：
\[
-
\frac{1}{\mu_{0}}F^{\mu\nu}\partial_{\mu}\delta A_{\nu}
=
-
\frac{1}{\mu_{0}}\partial_{\mu}(F^{\mu\nu}\delta A_{\nu})
+
\frac{1}{\mu_{0}}(\partial_{\mu}F^{\mu\nu})\delta A_{\nu}.
\]
相互作用项的变分为
\[
\delta\mathcal{L}_{\mathrm{int}}
=
-j^{\nu}\delta A_{\nu}.
\]

\subsection{非齐次 Maxwell 方程的推出}

忽略边界项后，作用量变分为
\begin{equation}
	\delta S
	=
	\frac{1}{c}
	\int
	\left[
	\frac{1}{\mu_{0}}\partial_{\mu}F^{\mu\nu}
	-
	j^{\nu}
	\right]
	\delta A_{\nu}\,\dif^{4}x.
	\label{eq:8.53}
\end{equation}
由于 $\delta A_{\nu}$ 任意，得到
\begin{equation}
	\partial_{\mu}F^{\mu\nu}
	=
	\mu_{0}j^{\nu}.
	\label{eq:8.54}
\end{equation}
这正是非齐次 Maxwell 方程 \eqref{eq:8.36}。齐次 Maxwell 方程则不是独立变分得到的动力学方程，而是由 $F_{\mu\nu}=\partial_{\mu}A_{\nu}-\partial_{\nu}A_{\mu}$ 自动推出的几何恒等式。

\section{Lorenz 规范与势的波动方程}

由
\begin{equation}
	F^{\mu\nu}
	=
	\partial^{\mu}A^{\nu}-\partial^{\nu}A^{\mu},
	\label{eq:8.55}
\end{equation}
非齐次 Maxwell 方程 \eqref{eq:8.36} 可写为

\[
\partial_{\mu}(\partial^{\mu}A^{\nu}-\partial^{\nu}A^{\mu})
=
\mu_{0}j^{\nu}.
\]

展开得
\begin{equation}
	\dalembert A^{\nu}
	-
	\partial^{\nu}(\partial_{\mu}A^{\mu})
	=
	\mu_{0}j^{\nu}.
	\label{eq:8.57}
\end{equation}
利用规范自由度，可以选择 Lorenz 规范
\begin{equation}
	\partial_{\mu}A^{\mu}=0.
	\label{eq:8.58}
\end{equation}
在三维形式中，

\[
\partial_{\mu}A^{\mu}
=
\frac{1}{c^{2}}\frac{\partial\phi}{\partial t}
+
\bm{\nabla}\cdot\bm{A}.
\]

因此 Lorenz 规范为
\begin{equation}
	\frac{1}{c^{2}}\frac{\partial\phi}{\partial t}
	+
	\bm{\nabla}\cdot\bm{A}=0.
	\label{eq:8.60}
\end{equation}
在 Lorenz 规范下，势方程化为
\begin{equation}
	\dalembert A^{\nu}=\mu_{0}j^{\nu}.
	\label{eq:8.61}
\end{equation}

时间分量给出
\[
\dalembert\phi=\frac{\rho}{\varepsilon_{0}},
\]
空间分量给出
\[
\dalembert\bm{A}=\mu_{0}\bm{j}.
\]
在无源区域中，$j^{\mu}=0$，于是
\[
\dalembert A^{\mu}=0.
\]
这表明在真空中，电磁势满足传播速度为 $c$ 的波动方程。

\section{带电粒子与电磁场的耦合}

\subsection{带电粒子的协变作用量}

带电粒子在外电磁场中的作用量应当是洛伦兹标量。自由粒子的作用量为
\[
S_{\mathrm{p}}
=
-mc\int \dif s.
\]
电磁相互作用项由四维势与世界线微元构造：
\[
S_{\mathrm{int}}
=
-q\int A_{\mu}\dif x^{\mu}.
\]
因此总作用量为
\begin{equation}
	S
	=
	\int
	\left(
	-mc\,\dif s
	-
	qA_{\mu}\dif x^{\mu}
	\right).
	\label{eq:8.67}
\end{equation}
将耦合项展开，利用 $A_{0}=\phi/c$、$\dif x^{0}=c\dif t$ 和 $A_{i}\dif x^{i}=-\bm{A}\cdot\dif\bm{r}$，有

\[
A_{\mu}\dif x^{\mu}
=
\phi\,\dif t
-
\bm{A}\cdot\dif\bm{r}.
\]

于是总作用量可写成
\[
S
=
\int_{t_{1}}^{t_{2}}
\left[
-mc^{2}\sqrt{1-\frac{v^{2}}{c^{2}}}
+
q\bm{A}\cdot\bm{v}
-
q\phi
\right]\dif t.
\]
对应的拉格朗日量为
\begin{equation}
	L
	=
	-mc^{2}\sqrt{1-\frac{v^{2}}{c^{2}}}
	+
	q\bm{A}\cdot\bm{v}
	-
	q\phi.
	\label{eq:8.70}
\end{equation}

\subsection{正则动量与力学动量}

由 \eqref{eq:8.70} 定义正则动量：

\[
\bm{P}
=
\frac{\partial L}{\partial\bm{v}}.
\]

直接计算得到
\begin{equation}
	\bm{P}
	=
	\gamma m\bm{v}+q\bm{A}.
	\label{eq:8.72}
\end{equation}
记力学动量为

\[
\bm{p}=\gamma m\bm{v},
\]

则
\begin{equation}
	\bm{P}=\bm{p}+q\bm{A}.
	\label{eq:8.74}
\end{equation}
在电磁场中，正则动量与力学动量不同。正则动量含有势函数项，力学动量对应粒子的实际运动。

哈密顿量为

\[
H=\bm{v}\cdot\bm{P}-L.
\]

代入 \eqref{eq:8.70} 与 \eqref{eq:8.72}，得

\[
H=\gamma mc^{2}+q\phi.
\]

于是

\[
H-q\phi=\gamma mc^{2}.
\]

由相对论能量动量关系得到
\begin{equation}
	(H-q\phi)^{2}
	=
	(\bm{P}-q\bm{A})^{2}c^{2}+m^{2}c^{4}.
	\label{eq:8.78}
\end{equation}
这就是带电粒子在外电磁势中的相对论质量壳关系。


由此可见，电磁势进入粒子动力学时，总是以组合
\[
H-q\phi,
\qquad
\bm{P}-q\bm{A}
\]
出现。这个结构通常称为最小耦合。它把自由粒子的能量动量关系改写为外电磁势中的质量壳关系，后面的 Lorentz 力方程也可从同一个拉格朗日量推出。

\subsection{Lorentz 力的三维形式}

由拉格朗日量 \eqref{eq:8.70} 出发，欧拉--拉格朗日方程为
\begin{equation}
	\frac{\dif}{\dif t}
	\left(
	\frac{\partial L}{\partial\bm{v}}
	\right)
	=
	\frac{\partial L}{\partial\bm{r}}.
	\label{eq:8.83}
\end{equation}
左侧为
\[
\frac{\dif}{\dif t}
\left(
\frac{\partial L}{\partial\bm{v}}
\right)
=
\frac{\dif\bm{p}}{\dif t}
+
q\frac{\partial\bm{A}}{\partial t}
+
q(\bm{v}\cdot\bm{\nabla})\bm{A}.
\]
右侧为
\[
\frac{\partial L}{\partial\bm{r}}
=
q\bm{\nabla}(\bm{A}\cdot\bm{v})
-
q\bm{\nabla}\phi.
\]

利用恒等式
\[
\bm{\nabla}(\bm{A}\cdot\bm{v})
=
(\bm{v}\cdot\bm{\nabla})\bm{A}
+
\bm{v}\times(\bm{\nabla}\times\bm{A}),
\]
代入 \eqref{eq:8.83} 后抵消公共项，得到
\[
\frac{\dif\bm{p}}{\dif t}
=
q\left(
-\bm{\nabla}\phi
-
\frac{\partial\bm{A}}{\partial t}
\right)
+
q\bm{v}\times(\bm{\nabla}\times\bm{A}).
\]
利用 \eqref{eq:8.8}，可得 Lorentz 力公式
\begin{equation}
	\frac{\dif\bm{p}}{\dif t}
	=
	q(\bm{E}+\bm{v}\times\bm{B}).
	\label{eq:8.88}
\end{equation}
磁力项不做功，因为 $(\bm{v}\times\bm{B})\cdot\bm{v}=0$。粒子能量变化率为
\begin{equation}
	\frac{\dif E}{\dif t}=q\bm{E}\cdot\bm{v}.
	\label{eq:8.89}
\end{equation}

\subsection{Lorentz 力的四维协变形式}

从协变作用量 \eqref{eq:8.67} 出发，对世界线作变分。自由粒子部分给出
\[
\delta S_{\mathrm{p}}
=
\int
\frac{\dif p_{\mu}}{\dif\tau}\delta x^{\mu}\,\dif\tau.
\]
相互作用项为
\[
S_{\mathrm{int}}
=
-q\int A_{\mu}\dif x^{\mu}.
\]
其变分经分部积分并利用端点固定条件后得到
\[
\delta S_{\mathrm{int}}
=
q\int F_{\mu\nu}u^{\nu}\delta x^{\mu}\,\dif\tau.
\]
总变分为
\[
\delta S
=
\int
\left[
\frac{\dif p_{\mu}}{\dif\tau}
+
qF_{\mu\nu}u^{\nu}
\right]
\delta x^{\mu}\,\dif\tau.
\]

由 $\delta x^{\mu}$ 任意，得
\begin{equation}
	\frac{\dif p_{\mu}}{\dif\tau}
	=
	-qF_{\mu\nu}u^{\nu}.
	\label{eq:8.94}
\end{equation}
升高指标后，得到最常用的协变 Lorentz 力方程：
\begin{equation}
	\frac{\dif p^{\mu}}{\dif\tau}
	=
	qF^{\mu}{}_{\nu}u^{\nu}.
	\label{eq:8.95}
\end{equation}
取空间分量可恢复三维 Lorentz 力 \eqref{eq:8.88}。

\section{对偶场张量与微分形式写法}

定义四维 Levi-Civita 张量满足
\[
\varepsilon^{0123}=1.
\]

电磁场张量的对偶张量定义为
\begin{equation}
	\widetilde{F}^{\mu\nu}
	=
	\frac{1}{2}\varepsilon^{\mu\nu\rho\sigma}F_{\rho\sigma}.
	\label{eq:8.97}
\end{equation}
齐次 Maxwell 方程 \eqref{eq:8.34} 可等价写为

\[
\partial_{\mu}\widetilde{F}^{\mu\nu}=0.
\]
因此 Maxwell 方程可以压缩为
\begin{equation}
	\partial_{\mu}F^{\mu\nu}=\mu_{0}j^{\nu},
	\qquad
	\partial_{\mu}\widetilde{F}^{\mu\nu}=0.
	\label{eq:8.99}
\end{equation}

若使用微分形式语言，四维势是一形式
\[
A=A_{\mu}\dif x^{\mu},
\]
电磁场是二形式
\begin{equation}
	F=\dif A.
	\label{eq:8.101}
\end{equation}
由于外微分满足 $\dif^{2}=0$，有
\begin{equation}
	\dif F=0.
	\label{eq:8.102}
\end{equation}
这正是齐次 Maxwell 方程。非齐次 Maxwell 方程则可写成
\begin{equation}
	\dif *F=\mu_{0}*J,
	\label{eq:8.103}
\end{equation}
其中 $J$ 是电流一形式，$*J$ 是它的 Hodge 对偶三形式。微分形式写法不是本章正文推导的必要工具，但它清楚揭示了齐次方程的几何来源。

\section{电磁场的 Lorentz 不变量}

由于 $F_{\mu\nu}$ 是张量，可以通过缩并构造 Lorentz 不变量。第一个不变量为
\begin{equation}
	I_{1}=F_{\mu\nu}F^{\mu\nu}.
	\label{eq:8.104}
\end{equation}
代入场张量分量可得
\begin{equation}
	I_{1}
	=
	2\left(B^{2}-\frac{E^{2}}{c^{2}}\right).
	\label{eq:8.105}
\end{equation}
第二个不变量由对偶张量给出：
\begin{equation}
	I_{2}=F_{\mu\nu}\widetilde{F}^{\mu\nu}.
	\label{eq:8.106}
\end{equation}
展开得
\begin{equation}
	I_{2}
	=
	-\frac{4}{c}\bm{E}\cdot\bm{B}.
	\label{eq:8.107}
\end{equation}
这两个量是 Lorentz 标量，可用于对电磁场作参考系无关的分类。

若
\[
\bm{E}\cdot\bm{B}=0,
\]
则电场和磁场正交。此时由 $c^{2}B^{2}-E^{2}$ 的符号决定场的类型。若 $cB>E$，称为类磁场，可以找到参考系使 $\bm{E}'=0$。若 $E>cB$，称为类电场，可以找到参考系使 $\bm{B}'=0$。若
\[
E=cB,
\qquad
\bm{E}\cdot\bm{B}=0,
\]
则两个不变量均为零，称为零场。平面电磁波就是典型例子。

\section{匀速运动点电荷的四维势}

考虑点电荷 $q$ 沿 $x$ 轴正方向以速度 $V$ 匀速运动。在共动参考系 $K'$ 中，电荷静止在原点，势函数为 Coulomb 势：
\begin{equation}
	\phi'
	=
	\frac{1}{4\pi\varepsilon_{0}}
	\frac{q}{\sqrt{x'^{2}+y'^{2}+z'^{2}}},
	\qquad
	\bm{A}'=0.
	\label{eq:8.110}
\end{equation}
从共动系 $K'$ 变换回实验室系 $K$，有
\[
\phi=\gamma\phi',
\qquad
A_{x}=\gamma\frac{V}{c^{2}}\phi'.
\]
因此
\begin{equation}
	A_{x}=\frac{V}{c^{2}}\phi,
	\qquad
	\bm{A}=\frac{\bm{V}}{c^{2}}\phi.
	\label{eq:8.112}
\end{equation}
坐标变换为
\[
x'=\gamma(x-Vt),
\qquad
y'=y,
\qquad
z'=z.
\]
代入 \eqref{eq:8.110}，得到实验室系中的标量势
\[
\phi(x,y,z,t)
=
\frac{1}{4\pi\varepsilon_{0}}
\frac{q\gamma}
{\sqrt{\gamma^{2}(x-Vt)^{2}+y^{2}+z^{2}}}.
\]
也可以写成
\begin{equation}
	\phi(x,y,z,t)
	=
	\frac{1}{4\pi\varepsilon_{0}}
	\frac{q}
	{\sqrt{(x-Vt)^{2}+
			\left(1-\frac{V^{2}}{c^{2}}\right)(y^{2}+z^{2})}}.
	\label{eq:8.115}
\end{equation}
实验室系中的势分布不再具有球对称性，横向距离项带有因子 $1-V^{2}/c^{2}$，这反映了运动电荷场分布的相对论性变形。

\section{本章小结}

本章从四维势出发建立了电磁场的协变描述。电磁势统一为四维矢量
\[
A^{\mu}
=
\left(
\frac{\phi}{c},
\bm{A}
\right),
\qquad
A_{\mu}
=
\left(
\frac{\phi}{c},
-\bm{A}
\right).
\]
电磁场张量由四维势定义为
\[
F_{\mu\nu}
=
\partial_{\mu}A_{\nu}-\partial_{\nu}A_{\mu}.
\]
在本章约定下，
\[
F_{0i}=\frac{E_{i}}{c},
\qquad
F_{ij}=-\varepsilon_{ijk}B_{k}.
\]
规范变换
\[
A'_{\mu}=A_{\mu}-\partial_{\mu}f
\]
不改变场张量。

Maxwell 方程的协变形式为
\[
\partial_{\mu}F^{\mu\nu}=\mu_{0}j^{\nu},
\]
以及
\[
\partial_{\lambda}F_{\mu\nu}
+
\partial_{\mu}F_{\nu\lambda}
+
\partial_{\nu}F_{\lambda\mu}=0.
\]
前者可由电磁场作用量对 $A_{\mu}$ 的变分推出，后者来自
$F_{\mu\nu}$ 由势函数定义这一几何结构。将这两个四维方程按时间分量和空间分量展开，即得到经典 Maxwell 方程组。
\[
	\left\{
	\begin{aligned}
		\bm{\nabla}\cdot\bm{E}
		&=
		\frac{\rho}{\varepsilon_{0}}, \\[6pt]
		\bm{\nabla}\cdot\bm{B}
		&=
		0, \\[6pt]
		\bm{\nabla}\times\bm{E}
		&=
		-\frac{\partial\bm{B}}{\partial t}, \\[6pt]
		\bm{\nabla}\times\bm{B}
		&=
		\mu_{0}\bm{j}
		+
		\frac{1}{c^{2}}\frac{\partial\bm{E}}{\partial t}.
	\end{aligned}
	\right.
\]
Lorenz 规范为
\[
\partial_{\mu}A^{\mu}=0.
\]

在该规范下，四维势满足
\[
\dalembert A^{\mu}=\mu_{0}j^{\mu}.
\]
带电粒子与电磁场的耦合由作用量项
\[
S_{\mathrm{int}}=-q\int A_{\mu}\dif x^{\mu}
\]
给出，并导出 Lorentz 力
\[
\frac{\dif\bm{p}}{\dif t}
=
q(\bm{E}+\bm{v}\times\bm{B}),
\qquad
\frac{\dif p^{\mu}}{\dif\tau}=qF^{\mu}{}_{\nu}u^{\nu}.
\]
电磁场的两个 Lorentz 不变量为
\[
F_{\mu\nu}F^{\mu\nu}
=
2\left(B^{2}-\frac{E^{2}}{c^{2}}\right),
\qquad
F_{\mu\nu}\widetilde{F}^{\mu\nu}
=
-\frac{4}{c}\bm{E}\cdot\bm{B}.
\]

本章建立了电磁场的四维表述：电场和磁场是反对称张量 $F_{\mu\nu}$ 的不同分量，Maxwell 方程、Lorentz 力和规范不变性都可以用协变形式表达。
